\documentclass[11pt]{article}

\usepackage[margin=25mm]{geometry}

\usepackage{amsmath,amssymb,fontspec,unicode-math}

\IfFontExistsTF{DejaVu Serif}{\setmainfont{DejaVu Serif}}{\setmainfont{Latin Modern Roman}}

\IfFontExistsTF{Latin Modern Math}{\setmathfont{Latin Modern Math}}{\setmathfont{DejaVu Math TeX Gyre}}

\usepackage{xurl,hyperref}

\setlength{\parindent}{0pt}\setlength{\parskip}{6pt}

\begin{document}

\title{Reinforcing Spatial Reasoning in Vision-Language Models with Interwoven Thinking and Visual Drawing}

\author{ELVIA preliminary AI analysis}\date{21 October 2026}\maketitle

Preliminary assessment: Minor.\par

Independent scientific review remains pending. Author adoption has not been requested.

The public reflection proxy loses step and line-operation information, and invalid view-change answers can receive credit through default zeros. The decoding-cost table omits generated-prefix interactions.

\section*{Review scope}

Latest-version, claim-focused review of the cited author-paper passages and applicable current author-release paths. The stated counterexamples and scope controls were checked analytically; no author program, old verifier, training run or benchmark was executed.

\section*{Problem 1: The reflection proxy does not retain cross-step operation identity}

Source: Sections 3.1–3.2, PDF pp.4–6, Eq.(5); Section 4.3, PDF p.9; verl/utils/reward\_score/gcot.py, format\_reward, lines 64–133

Assessment: supported; impact: minor

\subsection*{Source claim (paraphrase)}

The prose defines reflection across reasoning steps and includes box and line operations. `format\_reward` concatenates turns and compares only same-label box coordinates.

\subsection*{Analysis}

Two changed same-label boxes in one turn can therefore give a positive proxy without any second drawing step. Changed same-label lines across two turns are ignored. The printed formula itself omits an explicit distinct-step condition, so the first mismatch concerns its adjacent prose. Reflection is returned separately rather than simply added as RL reward. Author issue \#6 explains latest-same-label drawing merge as intentional behavior.

\subsection*{Finding and scope}

The public proxy does not reliably implement the cross-step all-operation description. Exact historical figure processing is unresolved.

\subsection*{Proposed correction}

Retain step, operation type, image identity, coordinates and label; require distinct steps and changed operation signatures for the stated syntactic proxy. This still does not prove semantic self-correction.

An explicit local correction is described; author adoption and a corresponding published erratum are not established.

\section*{Problem 2: Unknown motion fields silently become a valid zero vector}

Source: Appendix C.2, PDF p.26; released eval/evaluate.py, lines 24–73; SPAR annotation question 2863

Assessment: supported; impact: minor

\subsection*{Source claim (paraphrase)}

The view-change parser accepts arbitrary keys; recognized directional lookups default to zero before scoring.

\subsection*{Analysis}

Invalid text `bad:123` becomes the zero motion vector. Actual question \(2863\) has target \((1.6,0,0,-75,15)\), so two zero axes score perfectly and mean credit is \(0.4\). Across the \(400\) public view-change targets, zero-vector scores occur with counts \((129,117,91,50,13)\) at scores \((0,0.2,0.4,0.6,0.8)\), mean \(0.2505\). This is an artificial invalid-answer baseline, not any model’s observed prediction distribution. Explicit valid zero-motion answers may properly receive partial credit.

\subsection*{Finding and scope}

The evaluator needs an explicit invalid-answer policy. Published scores cannot be corrected without their raw predictions.

\subsection*{Proposed correction}

Validate recognized axes, duplicate/opposing keys and finite values before vector conversion, then apply the same valid-vector metric. Document any intentionally permissive omission convention.

An explicit local correction is described; author adoption and a corresponding published erratum are not established.

\section*{Problem 3: The decoding table omits within-response attention}

Source: Appendix E.1, PDF pp.30–31, Eqs.(9)–(10), Table 6; measured Table 7, PDF p.33

Assessment: supported; impact: minor

\subsection*{Source claim (paraphrase)}

Table 6 gives prompt-only \(O(MN)\) work for a dense cached decoder with input length \(M\) and output length \(N\).

\subsection*{Analysis}

Generated query \(j\) attends to \(M+j\) previous keys. Summing gives \(MN+N(N-1)/2\). At \(M=1,N=128\), the count is \(8256\), not \(128\); their ratio grows with output length. KV caching avoids recomputing old keys but not attention to them.

\subsection*{Finding and scope}

The unrestricted two-variable expression omits a term under standard dense attention. The measured runtime table and other sparse/recurrent mechanisms are not refuted.

\subsection*{Proposed correction}

For each generated block, include \(P_kB_k+B_k(B_k-1)/2\), where \(P_k\) is existing context and \(B_k\) the new block length. State a prompt-dominated approximation if that restricted regime is intended.

An explicit local correction is described; author adoption and a corresponding published erratum are not established.

\section*{Limits}

Preliminary AI assessment; independent human scientific adjudication remains pending.\par

An analytic counterexample or a release discrepancy does not by itself remeasure the paper’s reported experiments.\par

\section*{Source identity}

NeurIPS 2025 official proceedings PDF; compared with current arXiv:2506.09965v2

\url{https://papers.nips.cc/paper_files/paper/2025/file/d2f2ca2c305e1a9c29bf00a676f0a500-Paper-Conference.pdf}

Reviewed PDF SHA-256: \nolinkurl{152985272cefd6f66c283f58f1feb593229d5133aa859c78a553691fd11b1744}

The selected author/publisher source is the same work as the original case, matched by exact title, author identity and official version channel. The exact selected PDF SHA-256 was recomputed in the latest-version pass for Issue \#27. The PDF is a retained author/publisher copy; its original retrieval time is stated only when supported by matching acquisition metadata. File modification time was not substituted.

\end{document}
